\section{Global Inference y Provider Portfolio}

Cursor ejecuta modelos propios y, al mismo tiempo, llama a frontier providers. Autocomplete es extremadamente sensible a la latency: requiere GPU replicas en todo el mundo y un overload behavior definido. En las agent requests complejas importan más la calidad, el context y el tool use, y además dependen de los rate limits y de los cambios de versión de cada provider. Esta sección coloca el self-hosted serving y el multi-provider routing en un mismo plano de control de capacity/evaluation.

\subsection{Region, queue y fallback del self-hosted model}

Esta sección responde primero por qué autocomplete no puede desplegarse en un solo clúster centralizado. Al leer la figura, siga el recorrido por user region, router, GPU region y fallback: la distancia, la queue depth, las model replicas y la KV cache influyen en el p99. Cuando la region local se satura, el sistema debe elegir entre la latency remota, un smaller model y la graceful degradation.

\lecturefigure{11-global-inference.png}{La self-hosted inference requiere diseñar la latencia de cola interactiva con enrutamiento por region, capacidad y degradación.}{Subtítulos locales 03:20--04:20, 38:10--41:40; reelaboración conceptual.}

\begin{knowledgebox}{Lectura de la figura: el capacity headroom no es desperdicio}
El interactive serving no puede llevar la utilization promedio al 100\%. El queueing crece de forma no lineal cerca de la saturación, por lo que se necesitan replica headroom, admission control, deadlines y preemption. Las tareas por lote/re-embedding de baja prioridad pueden ocupar los periodos de baja demanda, pero deben ceder los recursos cuando aumenta el tráfico en línea.
\end{knowledgebox}

\subsection{Frontier providers: rate limit, quality, cost y failover}

\term{rate limit} es el límite que el provider impone a la concurrencia de requests/tokens o a una ventana de tiempo; \term{provider routing} elige el modelo según la tarea, la calidad, el estado de salud, la quota, el costo y la policy. Un round-robin simple ignora las diferencias de capacidad entre modelos y sus interfaces de context/tool; un portfolio realmente utilizable necesita normalized requests, versioned evals y un fallback específico para cada provider.

\lecturefigure{12-provider-routing.png}{La multi-provider inference requiere task contract, provider state, router y evidencia continua.}{Subtítulos locales 38:10--41:40; CursorBench y materiales posteriores sobre Cursor Router.}

\begin{knowledgebox}{Lectura de la figura: qué le falta a la offline eval y qué al online outcome}
El offline benchmark es reproducible y adecuado para pruebas de regresión, pero no refleja por completo los repos reales, la latency ni la user acceptance; la online metric se acerca al resultado del producto, pero la afectan la selección, la UI y el traffic mix. El router debe usar ambas y registrar el model/version en cada request y en cada incident.
\end{knowledgebox}

\begin{importantbox}{El routing es una optimización con restricciones}
Para una task $x$ y un model/provider candidato $m$, puede abstraerse como
\[
m^*=\arg\max_m Q(m,x)-\alpha C(m)-\beta L(m)
\]
sujeto a restricciones de privacy, availability, context, tool y quota. Los pesos $\alpha,\beta$ cambian según el plan, el tenant y la ruta de interacción, por lo que no conviene mantener un único «mejor modelo» global.
\end{importantbox}

\teachervoice{El realismo de la clase está en que los clientes de crecimiento rápido llaman una y otra vez para pedir más capacity, mientras la empresa prepara varios providers y GPU propias. La resilience no consiste en conectar un segundo proveedor de manera improvisada cuando el primero falla, sino en mantener en todo momento interfaces, evals, quotas y políticas de datos que permitan el cambio.}

\subsection{Resumen de la sección}

La global inference es a la vez un problema de queueing, capacity, evaluation y supply chain. Self-hosted y frontier provider no son opciones excluyentes, sino una combinación de alternativas con distinto grado de control, costo y modos de falla.
